% GENERATED FILE. Edit canonical lessons, then run npm run generate:books.
% canonical-source-hash: fnv1a64:8ca2de0ebe5fd2d0
% canonical-lessons: SA-C505-mrd, SA-C505-mrnmayah, SA-C505-mrdnati, SA-C505-apakah, SA-C505-akarah

\chapter{Clay and Fire}
\label{ch:sa-clay-and-fire}
\hlchaptermodality{505}

\begin{chapteropening}
\textbf{By the end of this chapter:} \emph{I can talk about how a potter works.}

Complete the last lesson of chapter 505: I can talk about how a potter works.
\end{chapteropening}

\section[\texorpdfstring{\sk{मृद्} (mṛd)}{मृद् (mṛd)}]{\sk{मृद्} (mṛd) — clay}

\label{lesson:SA-C505-mrd}



\begin{quote}
Before the new one: say the Sanskrit for lives by, then the Sanskrit for skilled.
\end{quote}

\subsection*{You'll want to know: \sk{मृद्}}
\textbf{\sk{मृद्}} — \emph{mṛd} — \textquotedblleft{}clay\textquotedblright{}.

\begin{quote}
\sk{कुम्भकारः} \sk{नद्याः} \sk{तीरात्} \sk{मृदम्} \sk{आनयति।}
\end{quote}

\begin{quote}
\emph{kumbhakāraḥ nadyāḥ tīrāt mṛdam ānayati.}
\end{quote}

\begin{quote}
The potter brings clay from the riverbank.
\end{quote}

A feminine noun.

\begin{cousinweb}
\textbf{\sk{मृद्}} (\emph{mṛd}) is the word MW records for earth, soil and clay, set beside the root √mṛd, \textquotedblleft{}to press, to crush\textquotedblright{}; MW gives no derivation, so its origin is uncertain (MW).
\end{cousinweb}

\begin{grammarlens}[title={What you've built}]
One more word for describing what you lived and felt.
\end{grammarlens}

\subsection*{Your turn}
\begin{itemize}
\raggedright
  \item \emph{Say it:} \emph{mṛd}
  \item \emph{Say it:} \emph{mṛd}, once more
  \item \emph{Say it:} \emph{kumbhakāraḥ nadyāḥ tīrāt mṛdam ānayati.} — The potter brings clay from the riverbank.
  \item \emph{Recall:} say the Sanskrit for lives by, then the Sanskrit for skilled, then say \emph{mṛd} again
\end{itemize}

\begin{tcolorbox}[breakable,colback=teal!4,colframe=teal!35!black,title={Before you move on}]
What does \sk{मृद्} mean? (\textquotedblleft{}Clay\textquotedblright{}.) Say it once more.
\end{tcolorbox}
\section[\texorpdfstring{\sk{मृन्मयः} (mṛnmayaḥ)}{मृन्मयः (mṛnmayaḥ)}]{\sk{मृन्मयः} (mṛnmayaḥ) — earthen}

\label{lesson:SA-C505-mrnmayah}



\begin{quote}
Before the new one: say the Sanskrit for skilled, then the Sanskrit for clay.
\end{quote}

\subsection*{You'll want to know: \sk{मृन्मयः}}
\textbf{\sk{मृन्मयः}} — \emph{mṛnmayaḥ} — \textquotedblleft{}earthen\textquotedblright{}.

\begin{quote}
\sk{कुम्भकारः} \sk{मृन्मयानि} \sk{पात्राणि} \sk{रचयति।}
\end{quote}

\begin{quote}
\emph{kumbhakāraḥ mṛnmayāni pātrāṇi racayati.}
\end{quote}

\begin{quote}
The potter makes earthen pots.
\end{quote}

\begin{cousinweb}
\textbf{\sk{मृन्मयः}} (\emph{mṛnmayaḥ}) is built on \emph{mṛd}, \textquotedblleft{}clay\textquotedblright{}, with \emph{-maya}, \textquotedblleft{}made of\textquotedblright{}; MW gives \emph{mṛnmaya}, so spelt (MW calls \emph{mṛṇmaya} a wrong reading), as made of earth or clay, earthen (MW).
\end{cousinweb}

\begin{grammarlens}[title={What you've built}]
One more word for describing what you lived and felt.
\end{grammarlens}

\subsection*{Your turn}
\begin{itemize}
\raggedright
  \item \emph{Say it:} \emph{mṛnmayaḥ}
  \item \emph{Say it:} \emph{mṛnmayaḥ}, once more
  \item \emph{Say it:} \emph{kumbhakāraḥ mṛnmayāni pātrāṇi racayati.} — The potter makes earthen pots.
  \item \emph{Recall:} say the Sanskrit for skilled, then the Sanskrit for clay, then say \emph{mṛnmayaḥ} again
\end{itemize}

\begin{tcolorbox}[breakable,colback=teal!4,colframe=teal!35!black,title={Before you move on}]
What does \sk{मृन्मयः} mean? (\textquotedblleft{}Earthen\textquotedblright{}.) Say it once more.
\end{tcolorbox}
\section[\texorpdfstring{\sk{मृद्नाति} (mṛdnāti)}{मृद्नाति (mṛdnāti)}]{\sk{मृद्नाति} (mṛdnāti) — kneads}

\label{lesson:SA-C505-mrdnati}



\begin{quote}
Before the new one: say the Sanskrit for clay, then the Sanskrit for earthen.
\end{quote}

\subsection*{You'll want to know: \sk{मृद्नाति}}
\textbf{\sk{मृद्नाति}} — \emph{mṛdnāti} — \textquotedblleft{}kneads\textquotedblright{}.

\begin{quote}
\sk{कुम्भकारः} \sk{मृदं} \sk{मृद्नाति।}
\end{quote}

\begin{quote}
\emph{kumbhakāraḥ mṛdaṃ mṛdnāti.}
\end{quote}

\begin{quote}
The potter kneads the clay.
\end{quote}

Past (he …): \sk{सः} \sk{मृदं} \sk{मृदितवान्}, \emph{saḥ mṛdaṃ mṛditavān}.

\begin{cousinweb}
\textbf{\sk{मृद्नाति}} (\emph{mṛdnāti}) comes from the root √mṛd, \textquotedblleft{}to press, to squeeze, to crush\textquotedblright{}, whose present is \emph{mṛdnāti}; its past participle is \emph{mṛdita} (MW).
\end{cousinweb}

\begin{grammarlens}[title={What you've built}]
One more word for describing what you lived and felt.
\end{grammarlens}

\subsection*{Your turn}
\begin{itemize}
\raggedright
  \item \emph{Say it:} \emph{mṛdnāti}
  \item \emph{Say it:} \emph{mṛdnāti}, once more
  \item \emph{Say it:} \emph{kumbhakāraḥ mṛdaṃ mṛdnāti.} — The potter kneads the clay.
  \item \emph{Recall:} say the Sanskrit for clay, then the Sanskrit for earthen, then say \emph{mṛdnāti} again
\end{itemize}

\begin{tcolorbox}[breakable,colback=teal!4,colframe=teal!35!black,title={Before you move on}]
What does \sk{मृद्नाति} mean? (\textquotedblleft{}Kneads\textquotedblright{}.) Say it once more.
\end{tcolorbox}
\section[\texorpdfstring{\sk{आपाकः} (āpākaḥ)}{आपाकः (āpākaḥ)}]{\sk{आपाकः} (āpākaḥ) — a kiln}

\label{lesson:SA-C505-apakah}



\begin{quote}
Before the new one: say the Sanskrit for earthen, then the Sanskrit for kneads.
\end{quote}

\subsection*{You'll want to know: \sk{आपाकः}}
\textbf{\sk{आपाकः}} — \emph{āpākaḥ} — \textquotedblleft{}a kiln\textquotedblright{}.

\begin{quote}
\sk{कुम्भकारः} \sk{पात्राणि} \sk{आपाके} \sk{स्थापयति।}
\end{quote}

\begin{quote}
\emph{kumbhakāraḥ pātrāṇi āpāke sthāpayati.}
\end{quote}

\begin{quote}
The potter puts the pots in the kiln.
\end{quote}

A masculine noun.

\begin{cousinweb}
\textbf{\sk{आपाकः}} (\emph{āpākaḥ}) is built on \emph{ā} and the root √pac, \textquotedblleft{}to cook, to bake\textquotedblright{}; MW gives \emph{āpāka} as a baking-oven, a potter's kiln (MW).
\end{cousinweb}

\begin{grammarlens}[title={What you've built}]
One more word for describing what you lived and felt.
\end{grammarlens}

\subsection*{Your turn}
\begin{itemize}
\raggedright
  \item \emph{Say it:} \emph{āpākaḥ}
  \item \emph{Say it:} \emph{āpākaḥ}, once more
  \item \emph{Say it:} \emph{kumbhakāraḥ pātrāṇi āpāke sthāpayati.} — The potter puts the pots in the kiln.
  \item \emph{Recall:} say the Sanskrit for earthen, then the Sanskrit for kneads, then say \emph{āpākaḥ} again
\end{itemize}

\begin{tcolorbox}[breakable,colback=teal!4,colframe=teal!35!black,title={Before you move on}]
What does \sk{आपाकः} mean? (\textquotedblleft{}A kiln\textquotedblright{}.) Say it once more.
\end{tcolorbox}
\section[\texorpdfstring{\sk{आकारः} (ākāraḥ)}{आकारः (ākāraḥ)}]{\sk{आकारः} (ākāraḥ) — a shape}

\label{lesson:SA-C505-akarah}



\begin{quote}
Before the new one: say the Sanskrit for kneads, then the Sanskrit for a kiln.
\end{quote}

\subsection*{You'll want to know: \sk{आकारः}}
\textbf{\sk{आकारः}} — \emph{ākāraḥ} — \textquotedblleft{}a shape\textquotedblright{}.

\begin{quote}
\sk{एतस्य} \sk{पात्रस्य} \sk{आकारः} \sk{सुन्दरः} \sk{अस्ति।}
\end{quote}

\begin{quote}
\emph{etasya pātrasya ākāraḥ sundaraḥ asti.}
\end{quote}

\begin{quote}
This pot has a beautiful shape.
\end{quote}

A masculine noun.

\begin{cousinweb}
\textbf{\sk{आकारः}} (\emph{ākāraḥ}) is built on \emph{ā} and the root √kṛ, \textquotedblleft{}to make\textquotedblright{}; MW gives \emph{ākāra} as form, figure, shape (MW).
\end{cousinweb}

\begin{grammarlens}[title={What you've built}]
One more word for describing what you lived and felt.
\end{grammarlens}

\subsection*{Your turn}
\begin{itemize}
\raggedright
  \item \emph{Say it:} \emph{ākāraḥ}
  \item \emph{Say it:} \emph{ākāraḥ}, once more
  \item \emph{Say it:} \emph{etasya pātrasya ākāraḥ sundaraḥ asti.} — This pot has a beautiful shape.
  \item \emph{Recall:} say the Sanskrit for kneads, then the Sanskrit for a kiln, then say \emph{ākāraḥ} again
\end{itemize}

\begin{tcolorbox}[breakable,colback=teal!4,colframe=teal!35!black,title={Before you move on}]
What does \sk{आकारः} mean? (\textquotedblleft{}A shape\textquotedblright{}.) Say it once more.
\end{tcolorbox}
